For each model, we first performed hyperparameter tuning on the training and development sets. The CNN and LSTM were tuned separately using grid search over the parameter ranges shown in Table~\ref{tab:hyperparameters}. During tuning, each configuration was trained for 3 epochs, and the best configuration was selected based on macro-F1 on the development set. Early stopping with a patience of 3 epochs was enabled during this tuning stage, although in practice the short tuning runs meant that most configurations finished quickly.

After selecting the best hyperparameters for a given model and sequence length, we trained that model again using the full training portion of the data and evaluated it on the unchanged test set. The final training runs used the Adam optimiser, a batch size of 512, gradient clipping with a maximum norm of 5.0, and 15 epochs due to early stopping. The final models were evaluated on the test sets.

We report accuracy and macro-F1 as our main evaluation metrics, with macro-F1 used as the primary measure for model selection because it gives equal weight to all four classes. In addition, we generated classification reports, confusion matrices, learning curves, and files with misclassified test examples for later error analysis. The test set was only used for final evaluation after model selection.

As our ablation study, we varied the maximum input sequence length while keeping the rest of the hyperparameters fixed to respective base models' constants. We ran each model with maximum sequence lengths of 64, 128, and 256 tokens, resulting in six main experiment settings: CNN-64, CNN-128, CNN-256, LSTM-64, LSTM-128, and LSTM-256. These experiments were executed on the Hábrók cluster as a job array, with one GPU assigned to each run. This design allows us to compare the two architectures directly and also study how much performance depends on the amount of input context available to the model.
